\documentclass{resume}

\address{1300 S University Ave \\ 503-A \\ Ann Arbor MI, 48104}
\address{(517) 539-1532 \\ sfalde@umich.edu}

\begin{document}

% Education
\begin{rSection}{Education}	
	\begin{rSubsection}{University of Michigan, Ann Arbor}{August 2023 - Present}{Sophomore Pursuing B.S.E. in Data Science}{GPA - 3.242}
		\item Programming (EECS 183, 280): Developed fundamental programming skills and concepts as applied to C++. Gained experience working as both a programming team member and leader with many projects including text comprehension, image resizing, and a terminal-based game.
		\item Data Structures (EECS 280, 281): Studied data structure concepts including linked lists, maps, inheritance, functors, iterators, polymorphism, and recursion, as applied to C++. Worked both individually and on a team for many projects including matrix implementations, a euchre game, and a machine learning social media post classifier. 
		\item Linear Algebra (MATH 214): Explored linear algebra concepts including linear equations, linear transformations, subspaces, and orthogonality. Applied concepts to linear programming, error-correcting codes, and 3D computer graphical computation.
		\item Discrete Mathematics (EECS 203): Studied mathematical concepts as they relate to computation including algorithms, counting, probability, induction, recursion, and logic.  % ! Finish
	\end{rSubsection}
\end{rSection}

% Activities
\begin{rSection}{Activities}
	\begin{rSubsection}{Michigan Data Science Team (MDST)}{August 2024 - Present}{Project Team Member}{}
		\item Applied programming skills and concepts to Python for data collection, cleaning, and analysis
	\end{rSubsection}
\end{rSection}

% Employment
\begin{rSection}{Employment}

	% Referee
	% Groundskeeping

\end{rSection}

% Strengths
\begin{rSection}{Strengths}
	\begin{tabular}{@{} >{\bfseries}l @{\hspace{2em}} l @{}}
		Programming & C++, Python, LaTeX, SQLite, R Studio \\ 
		Python & Pandas, NumPy, Matplotlib, Sklearn \\
		Software & Microsoft Office, Google Suite \\
	\end{tabular}
\end{rSection}

\end{document}
